\input{preamble}

\begin{document}

\title{}
\maketitle

\section{Can Machines Prove Theorems? 
Two Case Studies in AI-Assisted Mathematics}
\label{section-machines-prove-theorems}

\noindent
Ashvin Swaminathan, Harvard University.
June 19, 2026. Centro Pi Seminar, IMPA.

\medskip\noindent
{\bf Abstract.} 
Recent language models can do more than chat
about mathematics: they can complete proofs,
generate formal verifications, and
occasionally settle open questions. I will
examine what this means for research
mathematics through two concrete case studies
from my own work. The first is an old problem
in classical invariant theory that sat
unfinished in my notes for nearly a decade
and was completed using Claude Code and
Codex, with the key lemmas formally verified
in Lean 4 via Harmonic's Aristotle prover.
The second is even more striking: with
minimal prompting, I got GPT-5.5 to prove the
cubic analogue of a result of Laga and
myself, showing that a positive proportion of
hyperelliptic curves of large genus have no
points over cubic number fields. In both
cases, the AI contributed genuine
mathematical ideas, and formal verification
was essential to establishing that the
resulting arguments were correct, catching
errors in the process. I will use these
examples to discuss where automated theorem
proving and large language models are
genuinely useful today, where they still
fail, and why formalization is becoming more
important, not less, as these systems grow
more capable.

\medskip\noindent
Formal proof := every phrase can be computer
checked.

Proof assistant :=
programming language precise enough
to encode each proof step.

Lean. Needs libraries. 
Mathlib if a library of functions and
deftinitions (formalizations)
of mathematical concepts.
Currently = 115L defs vs 232K thms.
It's a language to write proofs that
computer can check.
Proof is written as code;
trusted kernel checks it has no gaps.
If it compiles, it's 100% correct.
Created by Leonardo de Moura
(from PUC-Rio!).

Proof that addition is commutative:

theorem add_comm( a b : Nat) : a+b=b+a:= by
induction b with
|zero => simp
|succ n ih =>rw[Nat.add_succ,ih,Nat.succ_add]

\medskip\noindent
simp means simplifiy
n+1 is denoted succ

sorry := it's an axiom

Proving := the agentic loop converges.

Example: 14,500 lines of Lean, 25 files.
Make the verifiable proof + check examples.

Claude Code = creative
Codex = more formal

-> there is not a very big infrastructure
for Lean these days!

Idea: combine codex + Lean.
This is what an "automated prover" does,
e.g. Aritstotle: leverage LLM
to create proofs in Lean.
Also AxiomProver.


Need to search mathlib to know what is already
available! Probably there's more
number theory/combinatorics than
geometry/topology (as of today).

\end{document}
